\section{Exact evaluation of a marked diagram}
\label{de:section}

The local polynomial now has a global interpretation. A variable belongs
to a face of the diagram, so all its occurrences at different boxes
must be identified before a coefficient is taken. We prove that this
shared coefficient is the reduced tensor evaluation, including its
single marked cap. No positivity or limiting estimate enters the proof.

\subsection{The ordered diagram and the coefficient theorem}

Throughout this section \(N\ge3\), \(d=N-1\),
\(q=e^{\pi i/N}\), \(Q=q^2\), and \(x_0=q^d\).
Start with a finite connected triangulation of the oriented sphere and
a perfect matching of its cubic dual. We allow a triangular
cellulation, but require that the two triangles at each matching edge
form an embedded quadrilateral disk with four distinct vertices.
There are \(2R\) triangles, \(R\) matching edges and \(R+2\) vertices.
The vertices will be called \emph{white faces}, because they index the
complementary regions in the diagram below.

Contract a small disk around each matched dual edge and its endpoints.
The disks are disjoint, and the resulting embedded graph \(\Gamma\)
is connected and four-valent, with \(R\) vertices and \(2R\) edges.
Choose an ordered parametrization of every disk. Its four surrounding
white faces have cyclic order
\[
                         (x,b,y,c);
\]
the matching diagonal is \(xy\). The coherently oriented triangles
can be written \((x,b,y)\) and \((y,c,x)\). The polynomial of this box
is \(H_f(x,b,c,y)\): \(b\) is its input middle face and \(c\) its
output middle face.

Here is the terminal convention. Write the four terminals as
\(C_0,C_1,C_2,C_3\) in cyclic order, with input terminals
\(C_1,C_2\) and output terminals \(C_0,C_3\). The matching halves
group \((C_1,C_2)\) and \((C_3,C_0)\); the untwisted physical
strands pair \((C_0,C_1)\) and \((C_2,C_3)\). These are different
pairings. When side indices \(0,1,2,3\) are used for a quadrilateral,
as in an annular braid drawing, they denote
\[
                 (C_1,C_2,C_3,C_0),
\]
respectively. In those indices the physical pairs are \((0,3)\)
and \((1,2)\), and the matching halves are \((0,1)\) and \((2,3)\).

Draw the two untwisted strands in each disk and join them along the
uncontracted dual edges. The outside arcs are disjoint. They can be
made Morse with respect to height, keeping the disk parametrizations
fixed, so that each maximum or minimum is an ordinary cap or cup.
Place the specified two-strand operator in each disk. A crossing
power uses the checked crossing of Proposition~\ref{fb:crossing-polynomial};
an arbitrary polynomial in the Casimir is also allowed. The local
frames are part of the data and are not reset when an arc is bent.

This construction is a literal planar tensor diagram. If crossing
circles are desired, place one perpendicular circle around the two
strands in each disjoint three-ball about a disk. Its crossing disk
meets precisely those two strands. Thus the corresponding augmented
link is embedded in \(S^3\); no geometric realization assumption is
needed for the finite evaluation. An arbitrary polynomial coupon
itself need not be a braid or define a knot.

\begin{figure}[ht]
 \centering
 \begin{tikzpicture}[x=1cm,y=1cm,>=Latex,
  face/.style={text=blue!55!black,font=\small},
  head/.style={font=\scriptsize},
  incidence/.style={blue!45!black,densely dashed,line width=.6pt},
  strand/.style={black,line width=.8pt}]
  % White vertices record the dual faces; black curves are the two strands.
  \begin{scope}
    \coordinate (x) at (-1.9,0);
    \coordinate (b) at (0,1.8);
    \coordinate (y) at (1.9,0);
    \coordinate (c) at (0,-1.8);
    \draw[incidence] (x)--(b)--(y)--(c)--cycle;
    \draw[incidence] (x)--(y);
    \foreach \p in {x,b,y,c}{\fill[blue!55!black] (\p) circle(1.5pt);}
    \node[face,left] at (x) {$x$};
    \node[face,above] at (b) {$b$};
    \node[face,right] at (y) {$y$};
    \node[face,below] at (c) {$c$};
    \draw[strand] (-.95,1.22)--(-.95,.4);
    \draw[strand] (.95,1.22)--(.95,.4);
    \draw[strand] (-.95,-.4)--(-.95,-1.22);
    \draw[strand] (.95,-.4)--(.95,-1.22);
    \draw[fill=white,rounded corners=2pt,line width=.7pt]
       (-1.3,-.4) rectangle (1.3,.4);
    \node at (0,0) {$f(\Omega)$};
    \node[head,left] at (-.95,1.25) {$i$};
    \node[head,right] at (.95,1.25) {$j$};
    \node[head,left] at (-.95,-1.25) {$r$};
    \node[head,right] at (.95,-1.25) {$s$};
    \node[head,left] at (-1,.88) {$C_1$};
    \node[head,right] at (1,.88) {$C_2$};
    \node[head,left] at (-1,-.88) {$C_0$};
    \node[head,right] at (1,-.88) {$C_3$};
    \node[font=\small] at (0,-2.45) {$i+j=r+s$};
  \end{scope}
  \begin{scope}[shift={(6.05,0)}]
    \draw[strand] (0,1.05)--(0,-.08);
    \node[head,above] at (0,1.05) {$i$};
    \node[face] at (-1.12,.48) {$L_x$};
    \node[face] at (1.12,.48) {$L_b$};
    \node[font=\small,fill=white,inner sep=3pt] at (0,-.55)
      {$i-\dfrac d2=\lambda_x-\lambda_b$};
    \node[font=\small,align=center] at (0,-1.55)
      {$L_f=\exp(2\pi\mathrm i\lambda_f/N)$\\[5pt]
       $Q^i=q^dL_x/L_b$};
    \node[font=\small,align=center] at (0,-2.45)
      {Centered head and face potentials};
  \end{scope}
\end{tikzpicture}

 \caption{The ordered box has cyclic face order \((x,b,y,c)\) and
 polynomial \(H_f(x,b,c,y)\). Its two inputs are \(C_1,C_2\), and
 its outputs are \(C_0,C_3\). The matching diagonal \(xy\) is not a
 physical strand. On a supported state the displayed head satisfies
 \(i-d/2=\lambda_x-\lambda_b\); this fixes the sign of the face
 potential. The alternative side numbering \(0,1,2,3\) refers to
 \(C_1,C_2,C_3,C_0\).}
 \label{de:figure}
\end{figure}

A main edge means one of the edges outside the box disks. Choose one
as a marked edge and call its adjacent faces \(a,b\). They are
distinct. Indeed a connected four-valent graph has no bridge:
if removing a single edge separated a set of vertices, the sum of
their degrees would be twice the number of internal edges plus one,
contrary to its even parity. An edge incident to only one face of
a spherical graph is a bridge.

The matrices defining the turns and mark are given below. Let
\(A_{\rm raw}(f_1,\ldots,f_R)\) be their contraction, with exactly
one marked cap and with the full operators \(f_e(\Omega)\) in the
recorded boxes. Each \(f_e\) is a finite polynomial; its coefficients
may depend on \(N\) and on other parameters. In particular it may
already include a signed crossing power through
Proposition~\ref{fb:crossing-polynomial}. Set
\[
 \mathcal P(L)=\prod_{e=1}^R
 H_{f_e}(L_{x_e},L_{b_e},L_{c_e},L_{y_e}).
\]
Every local factor is homogeneous of degree \(d\), so
\(\deg\mathcal P=dR\).

\begin{theorem}[The exact marked face coefficient]
\label{de:coefficient-theorem}
For the stated ordered diagram and the full polynomial coupons,
\begin{equation}\label{de:coefficient}
 A_{\rm raw}(f_1,\ldots,f_R)
 =\left[\prod_{f\ne a,b}L_f^d\right]
       \left.\mathcal P(L)\right|_{L_a=L_b=0}.
\end{equation}
This is an identity at each level, with no omitted power of \(N\).
The reduced scalar is independent of the chosen main edge when the
cut and its duality frames are transported as below. The theorem
does not assert this independence for separate projections of a
local polynomial onto selected monomials.
\end{theorem}

We prove the theorem in several finite steps. The first concerns
the reduced cut; the remaining steps turn its state sum into the
displayed coefficient.

\subsection{Duality, turns and the unrestricted mark}

For a generic nonzero parameter \(z\) which is not a root of unity,
let \(V_L\) have basis \(v_0,\ldots,v_L\) and operators
\[
 Kv_i=z^{L-2i}v_i,\qquad
 Fv_i=v_{i+1},\qquad
 Ev_i=[i]_z[L-i+1]_zv_{i-1}.
\]
Endpoint vectors are zero. The coproducts are those used in
Section~\ref{fb:crossing}, and the antipodes are
\[
 S(E)=-EK^{-1},\qquad S(F)=-KF,\qquad S(K)=K^{-1}.
\]
The dual action is obtained by transposing these antipode matrices.
Put
\begin{equation}\label{de:J}
 J_zv_i=c_i(z)v_{L-i}^*,\qquad
 c_i(z)=(-1)^iz^{i(i-L-1)}.
\end{equation}
This is an intertwining map from \(V_L\) to its dual. For \(F\),
the required ratio is
\[
 c_{i+1}/c_i=-z^{-L+2i};
\]
the \(K\) equation reverses the weights, and the \(E\) equation then
uses the same ratio and the displayed quantum integers. These
checks also show that \(J_z\) and its inverse are regular near our
root when \(L=d\).

The integer-weight operator \(H\) satisfies \(K=z^H\). Since \(E\)
and \(F\) change weights by two, \((-1)^H\) is central and
group-like. Consequently
\[
              \widetilde g=(-1)^HK
\]
implements \(S^2\) on every tensor product and dual used here.
In the given self-dual frame it is exactly the double-dual map:
\[
 J_z^{-t}J_z=\widetilde g,\qquad
 \frac{c_i(z)}{c_{L-i}(z)}=(-1)^Lz^{L-2i}.
\]
At \(z=q\), every integral weight satisfies \(q^{-NH}=(-1)^H\).
Thus the generic pivot specializes to the physical root pivot
\begin{equation}\label{de:pivot}
                       g=K^{1-N}.
\end{equation}
This specialization uses a regular generic pivot; it does not
assert that \(K^{1-N}\) implements \(S^2\) for a generic \(z\).

Take \(L=d\) and pass to the balanced basis
\(e_i=v_i/(Q;Q)_i\) at the root. The product identity
\begin{equation}\label{de:factorial-cancellation}
 (Q;Q)_i(Q;Q)_{d-i}
       =N(-1)^iQ^{i(i+1)/2}
\end{equation}
follows from \((Q;Q)_d=N\) and
\(1-Q^{-j}=-Q^{-j}(1-Q^j)\).
Since \(c_i(q)=q^{i^2}\), the factorials in \(J_q\) therefore
cancel exactly. If the first tensor factor below is a dual factor,
denote it by a superscript \(\vee\). The four ordered turns are
\[
\begin{aligned}
 \operatorname{ev}_L(e_a^\vee\otimes e_b)
      &=N^{-1}x_0^a\,\mathbf1_{a+b=d},&
 \operatorname{coev}_L(1)
      &=\sum_a Nx_0^{-a}e_{d-a}\otimes e_a^\vee,\\
 \operatorname{ev}_R(e_b\otimes e_a^\vee)
      &=N^{-1}x_0^ag_b\,\mathbf1_{a+b=d},&
 \operatorname{coev}_R(1)
      &=\sum_a Nx_0^{-a}g_{d-a}^{-1}
                       e_a^\vee\otimes e_{d-a},
\end{aligned}
\]
where
\[
                       g_j=q^{-d^2-2j}.
\]
Contraction of a turn with the opposite turn has a single matching
index; its \(x_0\) and \(g\) factors cancel. This proves all four
snake identities. Closing an endomorphism by the corresponding
ordinary turns gives the pivotal trace \(\operatorname{Tr}(gA)\).

In literal left-to-right physical head order \((i,j)\), both cups
and both caps reduce to
\begin{equation}\label{de:root-turns}
 p_{ij}=Nq^{-d^2+di}\mathbf1_{i+j=d},
 \qquad
 \ell_{ij}=N^{-1}q^{di}\mathbf1_{i+j=d}.
\end{equation}
For example the right-cap ratio to the stated cap is
\(q^{-2Ni}=1\), and the right-cup ratio is
\(q^{2N(d-i)}=1\). These are equalities of the actual turn matrices,
not independent rescalings of the local frames.

Define the rank-one projector \(e=\psi\phi\) by
\(\psi_i=1\), \(\phi_i=1/N\), so \(e^2=e\) and
\(\operatorname{Tr}e=1\). At a right cut insert \(g^{-1}e\)
on the primal input. For every matrix \(C\),
\[
                 \operatorname{Tr}(g\,g^{-1}eC)
                           =\operatorname{Tr}(eC).
\]
At a left cut use the corresponding defect \(ge\). In physical
left-to-right order the resulting marked rows are
\begin{equation}\label{de:unrestricted-mark}
 m_R(i,j)=\frac{x_0^j}{N^2},\qquad
 m_L(i,j)=\frac{x_0^i g_{d-i}}{N^2}.
\end{equation}
For the right row, the only summand in
\(\operatorname{ev}_R(g^{-1}e\,e_i\otimes e_j^\vee)\)
has primal index \(d-j\), and its pivot times defect entry is
\(1/N\). This gives the first formula; the left formula follows
by the same single-index contraction. Neither marked row has a
support restriction.

Each complete marked cup-cap contraction is one. In a whole
closed diagram with one marked cap, conservation forces the
marked heads to satisfy \(i+j=d\): each box preserves the true
integer sum of its two indices, every cup supplies \(d\), and
every other cap consumes \(d\). On this forced support both
rows in \eqref{de:unrestricted-mark} equal
\begin{equation}\label{de:mark-on-support}
                         m_{ij}=N^{-2}q^{dj}.
\end{equation}
We shall use \eqref{de:mark-on-support} only after this global
conservation argument. Restricting the marked row in advance
would be a different tensor.

\subsection{Why the chosen reduced cut can be exposed}

We next justify moving the marked edge to an exposed top cap
while keeping every coupon disk and its ordered frame fixed.
This is the only place where a closed spherical reading of the
diagram is compared with a different planar reading.

For the generic pivot \(\widetilde g\), the two traces agree on
every intertwiner in the tensor-generated category:
\begin{equation}\label{de:spherical-traces}
 \operatorname{Tr}(\widetilde gA)
       =\operatorname{Tr}(\widetilde g^{-1}A).
\end{equation}
Here is a direct proof including multiplicity spaces. First take
real \(z>0\), \(z\ne1\). The involution
\[
                  K^*=K,\qquad E^*=KF,\qquad F^*=EK^{-1}
\]
is compatible with the stated coproducts. On \(V_L\), choose
orthogonal basis vectors with squared norms
\[
 n_0=1,\qquad
 n_i=[i]_z[L-i+1]_z\,z^{-L+2i}n_{i-1}>0.
\]
These give precisely the displayed adjoint relations. Tensor
products are unitary for the tensor-product norm; duals are
unitarized by \eqref{de:J}. An invariant subspace therefore has
an invariant orthogonal complement.

Every irreducible summand is one of the \(V_L\). Indeed choose a
highest vector \(v\) of integral weight \(L\). The identity
\[
 EF^rv=[r]_z[L-r+1]_zF^{r-1}v
\]
shows that the first zero lowering occurs at \(r=L+1\); a
negative highest weight would have no zero lowering. Bounded
weights exclude that possibility. The preceding nonzero
lowerings have distinct weights and span the irreducible
summand. An intertwiner consequently has blocks
\(I_{V_L}\otimes A_L\). On each such block both sides of
\eqref{de:spherical-traces} equal
\[
                         (-1)^L[L+1]_z\operatorname{Tr}(A_L),
\]
because the weights \(L,L-2,\ldots,-L\) are symmetric.

For a fixed diagram all matrix entries are rational in a fixed
local square root of \(z\). Polynomial coupon coefficients are
held fixed while taking this generic lift. The identity on an
interval of positive real \(z\) proves the rational identity,
and all original matrices are regular at \(q\). For the
crossing this regularity was proved in
Proposition~\ref{fb:crossing-polynomial}; its inverse is regular
as well because its root value is invertible. Thus
\eqref{de:spherical-traces} specializes to the actual diagram.
No decomposition or projection onto the generic irreducible
summands is specialized.

Choose a point in face \(a\) and move it to the point at infinity
along a path outside all coupon disks. A thin strip around that
path gives a sphere isotopy which fixes the disks pointwise.
The outside arcs may be crossed by the moving infinity, but no
crossing or ribbon twist is introduced. Ordinary portions of the
motion are tensor interchange and the four snake identities.
When infinity passes an arc, the two planar readings close the
same bent endomorphism by its two pivotal traces. Equality
\eqref{de:spherical-traces} gives the same generic closed scalar.
Face \(a\) is now exterior.

The chosen edge borders that exterior face. Push a short subarc
along a narrow strip within the face to a point above the
remaining compact drawing. Its tip is a top cap; the strip
introduces no crossing or extra component. The inside face is
still \(b\). Perturb the other arcs to make their turns Morse.

Leave the new cap open. After the actual cup identification,
the remaining diagram is an endomorphism of \(V_d\) or its dual.
Generically it is a scalar \(a(z)I\): commuting with the distinct
\(K\) weights makes it diagonal, and commuting with the nonzero
adjacent \(F\) arrows makes its diagonal constant. The scalar is
regular at \(q\), since it is an entry of a regular open network.
Closing it ordinarily gives
\[
                         (-1)^d[N]_z\,a(z).
\]
This factor is the same for every exposed cut, even on a
different component, because every component has the same
color. It is nonzero generically, so equal closed scalars have
equal regular open scalars. Specialization gives their equality
at \(q\), without division by \([N]_q=0\).
The marked closure evaluates \(\operatorname{Tr}(e\,aI)=a\).
This proves the required independence of the reduced cut.

This argument does not commute the nonintertwining defect
\(g^{-1}e\) through the boxes. It first transports the cut with
its true cups and duality, and then identifies its reduced scalar.

\subsection{The affine face lattice and the true support}

We describe the state space before evaluating the turn factors.
For a vertical head separating a left face \(f\) from a right
face \(g\), use the convention
\begin{equation}\label{de:centered-flow}
 i-\frac d2=\lambda_f-\lambda_g,\qquad
 L_f=\exp(2\pi i\lambda_f/N).
\end{equation}
At the first input of a box this reads
\(Q^i=q^dL_x/L_b=-q^{-1}L_x/L_b\), exactly the first
clock in the local Fourier proof. At the other input the ratio
is \(q^dL_b/L_y\); the two output ratios use \(c\) in place of
\(b\). Reversing an arc replaces its head \(i\) by \(d-i\),
so its centered value changes sign.

A physical state on the unmarked diagram is a choice of canonical
heads \(0,\ldots,d\). At a turn the heads sum to \(d\); at every
box the two input heads have the same integer sum as the two
output heads. Centered heads thus define an antisymmetric flow
on the edges of \(\Gamma\), with zero divergence at each vertex.
Equivalently, they give a closed cochain on the dual face graph.

For clarity, the required integral fact holds over every
abelian group. A zero-divergence flow on a connected spherical
graph is the difference of face potentials, unique after one
potential is fixed. To prove it, integrate the dual cochain from
the fixed face along dual edge paths. Two paths differ by a
closed dual walk. Splitting repeated vertices reduces that walk
to simple cycles; on the sphere each such cycle bounds a union
of dual two-cells. Its cochain sum is the sum of the divergences
at the enclosed primal vertices, hence zero. This proves
path independence. Connectedness gives uniqueness. Integration
uses only addition and subtraction, so it introduces no finite
index or determinant factor.

For odd \(N\), the centered heads are integral. Fix
\(\lambda_a=0\); their potentials are then integers. For even
\(N\), twice every centered head is odd. The corresponding
parity cochain equals one on each main edge and zero on a
matching diagonal: the boundary of each triangle has two
main edges. It is closed, and the same spherical argument over
\(\mathbb F_2\) gives unique face parities \(t_f\), with
\[
 t_a=0,\qquad t_b=1.
\]
Equivalently, the dual quadrilateral at each box has four odd
edges and opposite vertices of the same parity. The exact
potentials have the form
\begin{equation}\label{de:affine-potentials}
                  \lambda_f=p_f+\frac{t_f}{2},
                  \qquad p_f\in\mathbb Z.
\end{equation}
For odd \(N\), use \(t_f=0\) throughout. These affine lattices are
nonempty: the potential \(t_f/2\) itself gives centered flow
\((t_f-t_g)/2\), and its heads lie in the physical range.
The difference of any two lifts is an integral face potential.

Reduce the \(p_f\) modulo \(N\), with \(p_a=0\). There are exactly
\(N^{R+1}\) face-grid points
\begin{equation}\label{de:face-grid}
 L_a=1,\qquad L_f=q^{t_f}Q^{p_f}\quad(f\ne a).
\end{equation}
At every such point take the canonical head representing
\(d/2+\lambda_f-\lambda_g\) modulo \(N\). At a turn the
congruence \(i+j=d\pmod N\) forces the equality \(i+j=d\):
neither \(d-N=-1\) nor \(d+N=2N-1\) lies in \([0,2d]\).

At a box the residue total need not determine the true integer
total. This is the place where physical support is essential.
The local Fourier proposition expresses the actual residue
block as a direct sum of its true-total blocks. Hence if the
canonical input and output totals differ, the matrix entry
vanishes, and \eqref{uf:physical-local-adapter} implies
\begin{equation}\label{de:unsupported-zero}
                    H_{f_e}(L_x,L_b,L_c,L_y)=0
\end{equation}
at that grid point. The prefactor there is nonzero.

If no box violates the true total, the canonical heads form the
closed centered flow just described. Integrating it gives unique
exact potentials with \(\lambda_a=0\), and they reduce to the
given grid point. Conversely every supported canonical state
gives one grid point by this integration. Additional zeros of
a coupon merely give zero weight on both sides. We have proved
a bijection between supported states and supported grid points,
and proved that all other grid points contribute zero.

The exact potentials of a supported state, rather than arbitrary
representatives \(0,\ldots,N-1\) for the \(p_f\), are used in
intermediate half-powers. In particular at a box
\[
                 r-u=\lambda_b-\lambda_c
\]
is the canonical integer displacement, since \(b,c\) have the
same parity. Thus
\begin{equation}\label{de:coherent-adapter}
 q^{-d(r-u)}c^{-d}=(bc)^{-d/2},\qquad
 f_e(\Omega)_{r,s;u,v}
        =\frac1N\frac{H_{f_e}(x,b,c,y)}{(bc)^{d/2}}.
\end{equation}
Here each fractional power means the exponential of the exact
potential in \eqref{de:centered-flow}. At even \(N\) one has
\(q^{d(p+N)}=-q^{dp}\); the proof has not asserted otherwise.
An unsupported grid point needs no such lift, because its
product already vanishes by \eqref{de:unsupported-zero}.

\subsection{Winding, the single cut factor and coefficient extraction}

Let \(M_f\) be the number of input and output middle occurrences
of face \(f\) at the boxes. Let \(T_f\) count outside turn
occurrences positively and inside turn occurrences negatively,
temporarily treating the marked cap as ordinary.
Then
\begin{equation}\label{de:winding}
 M_f-T_f=2\quad(f\ne a),\qquad M_a-T_a=-2.
\end{equation}
Indeed traverse each face boundary with the face on its left.
Resolve repeated boundary passages by a narrow regular
neighborhood. In the recorded local box resolution each middle
passage reverses the vertical tangent and contributes \(\pi\);
the two outer passages have net turn zero. An inside cup or cap
contributes \(\pi\), and its outside passage contributes
\(-\pi\). The intervening monotone arcs have zero net contribution
after deformation to vertical tangents. This also follows by
counting signed passages through a horizontal tangent direction.
The bounded-face boundary has winding \(+1\), and the exterior
boundary winding \(-1\), proving \eqref{de:winding}. Every repeated
occurrence of a face is counted, and no restriction on the
relative heights of different boxes was used.

At a cup or cap with outside face \(o\) and inside face \(z\),
the first head satisfies
\[
 q^{di}=q^{d^2/2}(L_o/L_z)^{d/2}.
\]
Thus an ordinary cup contributes \(Nq^{-d^2/2}\) times this
face character and a cap contributes \(N^{-1}q^{d^2/2}\) times
the same character. There are equally many cups and caps, so
their constant phases and powers of \(N\) cancel before one cap
is marked. At that exposed cap, \eqref{de:mark-on-support} gives
\begin{equation}\label{de:mark-ratio}
 \frac{m_{ij}}{\ell_{ij}}
      =N^{-1}q^{d(j-i)}
      =N^{-1}(L_b/L_a)^d.
\end{equation}
Equations \eqref{de:coherent-adapter} and \eqref{de:winding}
therefore turn the whole literal marked summand into
\begin{equation}\label{de:marked-summand}
 N^{-R-1}\,
       \frac{\mathcal P(L)}{\prod_{f\ne a,b}L_f^d}.
\end{equation}
Equivalently, with \(k\) cups, \(k-1\) ordinary caps and one
marked cap, the turn count is
\(N^kN^{-(k-1)}N^{-2}=N^{-1}\); each of the \(R\) local Fourier
inversions supplies its separate \(N^{-1}\).
This is the single reduced-cut normalization.

All fractional face powers have disappeared from
\eqref{de:marked-summand}. Changing any exact potential by a
multiple of \(N\) leaves this whole expression unchanged. By
the supported-state argument it can be summed over the entire
grid \eqref{de:face-grid}, not just over a chosen lift.

It remains to calculate this normalized grid average.
A monomial of \(\mathcal P\) has nonnegative face exponents \(n_f\)
with total \(dR\). Character orthogonality on each of the
\(R+1\) independent \(N\)-point grids gives the conditions
\[
 n_b\equiv0\pmod N,\qquad n_f\equiv d\pmod N
                       \quad(f\ne a,b).
\]
There are \(R\) ordinary faces, each with exponent at least \(d\).
Their minimum already exhausts the total degree \(dR\).
Consequently the only surviving exponents are
\begin{equation}\label{de:unique-target}
             n_a=n_b=0,\qquad n_f=d\quad(f\ne a,b).
\end{equation}
At this exponent the even-coset character is exactly one.
The \(N^{R+1}\) grids cancel the denominator in
\eqref{de:marked-summand}. This proves
Theorem~\ref{de:coefficient-theorem}, with both parities and the
entire normal data of every coupon retained.

\subsection{Insertion of signed powers after the shared extraction}

The preceding theorem is an equality for full polynomial
operators. It may therefore be combined with resolvent
reproduction before any estimate is taken. Define
\begin{equation}\label{de:shared-numerator}
 S_N(\boldsymbol\sigma)=
 \left[\prod_{f\ne a,b}L_f^d\right]
 \left.
 \prod_{e=1}^R H_{V_N(\sigma_e;\Omega)}
        (L_{x_e},L_{b_e},L_{c_e},L_{y_e})
 \right|_{L_a=L_b=0}.
\end{equation}
Each spectral degree is at most \(d\). The extraction and
multiplication in \eqref{de:shared-numerator} take place in the
same face variables as in the original diagram.

\begin{corollary}[Exact filling of the marked source]
\label{de:exact-filling}
For any fixed signed integer powers \(n_1,\ldots,n_R\), the
raw marked evaluation with \(B_q^{n_e}\) in box \(e\) is
\begin{equation}\label{de:tensor-filling}
 A_{\rm raw}(\boldsymbol n)
       =\left(\bigotimes_{e=1}^R\mathcal L_N^{n_e}\right)
                                        S_N(\boldsymbol\sigma).
\end{equation}
If \(C_N\) is any other polynomial with the same separate spectral
degree bound, then
\[
 \left|A_{\rm raw}(\boldsymbol n)
       -\left(\bigotimes_e\mathcal L_N^{n_e}\right)C_N\right|
 \le\prod_e(1+|n_e|/2)\,
                   \|S_N-C_N\|_{[-2,2]^R}.
\]
The joint-grid improvement of Theorem~\ref{fb:main} applies when
its vanishing hypothesis is satisfied.
\end{corollary}

\begin{proof}
Proposition~\ref{prop:resolvent-reproduction} gives
\(\mathcal L_N^{n_e}V_N(\sigma_e;\Omega)=
\widetilde f_{n_e}(\Omega)\) in each individual box.
Proposition~\ref{fb:crossing-polynomial} identifies the latter
with \(B_q^{n_e}\). The ordered tensor contraction is multilinear
in its local matrices, so these identities may be applied
separately at every box and then
Theorem~\ref{de:coefficient-theorem} gives
\eqref{de:tensor-filling}. No commutation between boxes is used.
Subtract the two full shared polynomials before applying
Theorem~\ref{fb:main}; this proves the norm assertion.
\end{proof}

The exact coefficient identity and the last corollary contain
no asymptotic remainder hypothesis. A strict bound for the
particular polynomial \(S_N-C_N\), and any nonzero limiting
contribution assigned to \(C_N\), remain separate analytic
questions.


\section{Comparison with the reduced Kashaev invariant}
\label{de:normalization-section}

We specify the invariant convention rather than leaving a phase or
a change of basis implicit. This comparison is finite and applies
at every \(N\ge2\). The global coefficient theorem above is used
only at \(N\ge3\).

Let \(\mathsf P\) exchange tensor factors. In this section set
\(R=\mathscr R_q\) and \(B=B_q=\mathsf P R\), using the
crossing defined in Section~\ref{fb:crossing}. Set
\[
 p_j=\prod_{r=1}^j(1-Q^r),\qquad
 \bar p_j=\prod_{r=1}^j(1-Q^{-r}),\qquad p_0=\bar p_0=1.
\]
For an integer \(a\), let \([a]_N\) be its representative in
\(\{0,\ldots,d\}\). The enhanced Kashaev crossing in the
column-vector convention is
\begin{equation}\label{de:K-crossing}
 \begin{aligned}
 (B_K)_{ce,ab}&=
 \frac{NQ^{1+c-b+(a-e)(c-b)}
       \mathbf1_{u+t<N}\mathbf1_{v+h<N}}
      {p_u\bar p_vp_t\bar p_h},\\
 (u,v,t,h)&=([b-a-1]_N,[a-c]_N,[c-e]_N,[e-b]_N).
 \end{aligned}
\end{equation}
Its enhancement is
\[
 \mu_K=-qS,\qquad \alpha_K=-q,\qquad \beta_K=1,
 \qquad Sv_j=v_{[j+1]_N}.
\]
These are the enhanced crossing and normalization of
Murakami--Murakami~\cite[\S\S1,4]{MurakamiMurakami2001}, with
their quantum parameter equal to \(Q\).
A negative braid generator uses the actual inverse of
\eqref{de:K-crossing}. The reduced value is obtained by cutting
one strand, taking the \(\mu_K\)-weighted traces on the others,
and multiplying by \(\alpha_K^{-w}\), where \(w\) is the
writhe of the braid word. Its unknot value is one.
We use the standard enhanced ribbon invariant defined by these
data; the calculation below proves the normalization comparison
rather than invoking its comparison theorem.

\begin{theorem}[Exact normalization of the marked coefficient]
\label{de:invariant-normalization}
In the enhanced braid conventions just specified, a link \(L\)
with \(c\) components satisfies
\begin{equation}\label{de:enhanced-comparison}
 F_N^g(L)=(-1)^{d(c-1)}K_N(L),
 \qquad g=K^{1-N},
\end{equation}
where \(F_N^g\) is the reduced \(g\)-pivotal evaluation normalized
by its positive curl. For a knot the sign is one.

For a diagram \(D\) of a knot \(K\) in
Theorem~\ref{de:coefficient-theorem}, with its boxes filled by
the specified signed crossing powers, the zero-framed reduced
invariant is
\begin{equation}\label{de:raw-invariant}
 K_N(K)=\zeta_{N,D}A_{\rm raw}(\boldsymbol n),
              \qquad |\zeta_{N,D}|=1.
\end{equation}
The character \(\zeta_{N,D}\) consists of the recorded integer
framing corrections. For a closed braid of writhe \(w\) in the
same convention, it is \(\theta^{-w}\), where
\(\theta=q^{-d^2/2}\).
Consequently the right side of \eqref{de:tensor-filling}, with
this same character, evaluates the original invariant.
\end{theorem}

\begin{proof}
We first compare the matrices, keeping all factorial changes
of basis. Define \(R^o=(\mathsf P B_K)^{\mathsf T}\) and the
Fourier matrix \(\mathsf F_{ij}=Q^{ij}/\sqrt N\). Put
\[
 A=(\mathsf F^{-1}\otimes\mathsf F^{-1})R^o
                         (\mathsf F\otimes\mathsf F).
\]
We claim that its only nonzero entries are
\begin{equation}\label{de:triangular-Fourier}
 A_{a-k,b+k;\,a,b}
   =Q^{1+(a-k)(b+1)}
      \frac{p_a p_{b+k}}{p_{a-k}p_b p_k},
 \qquad 0\le k\le\min(a,d-b).
\end{equation}
In particular the conservation is equality of integer sums.

Here is the finite calculation. For a nonzero entry of \(R^o\)
introduce the residues
\[
 u=[j-i-1]_N,\quad v=[i-l]_N,\quad
 t=[l-k]_N,\quad h=[k-j]_N.
\]
The two indicators say \(u+t<N\) and \(v+h<N\); their total
is congruent to \(-1\) modulo \(N\), so
\(u+v+t+h=d\). Conversely this composition and one free index
\(j\) determine the original entry:
\[
 i=j-1-u,\quad k=j+h,\quad l=j+h+t\pmod N,\qquad
 (R^o)_{ij,kl}
 =\frac{NQ^{1-(t+h)(u+h)}}{p_u\bar p_vp_t\bar p_h}.
\]
Summing the free index in the Fourier transform gives the
modular conservation delta. If
\[
 S_r(X)=\sum_{j=0}^r\frac{X^j}{p_j\bar p_{r-j}},
\]
the remaining sum is
\begin{equation}\label{de:Fourier-sum}
 A_{c,e;\,a,b}
 =\mathbf1_{a+b\equiv c+e\pmod N}
 \sum_{z=0}^d Q^{1+c-z^2+z(a+b)}
             S_{d-z}(Q^{c-z})S_z(Q^{z-a}).
\end{equation}
Use
\(\bar p_r=(-1)^rQ^{-r(r+1)/2}p_r\).
Comparing coefficients in
\((1-Q^r)S_r(X)=(X-Q^r)S_{r-1}(X)\) proves
\[
                    S_r(X)=\frac{\prod_{j=1}^r(X-Q^j)}{p_r}.
\]
The two factors in \eqref{de:Fourier-sum} now force
\(c\le z\le a\). With \(k=a-c\), the surviving sum is
\[
 Q^{1+c(b+1)}\frac{p_a}{p_c}
 \sum_{j=0}^k
   \frac{(-1)^jQ^{j(j-1)/2}Q^{(b+1)j}}{p_jp_{k-j}}
 =
 Q^{1+c(b+1)}\frac{p_a}{p_c}
                    \frac{(Q^{b+1};Q)_k}{p_k}.
\]
Here \((a;Q)_k=\prod_{r=0}^{k-1}(1-aQ^r)\), with empty
product one. The last equality is the same finite product identity after
reversing its coefficients. If \(b+k\ge N\), the numerator
contains \(1-Q^N\) and is zero. Otherwise
\((Q^{b+1};Q)_k=p_{b+k}/p_b\). This proves
\eqref{de:triangular-Fourier} with its endpoint support.

Define diagonal matrices and a scalar by
\[
 G_j=\frac{p_j}{(1-Q)^j}=q^{j(j-1)/2}[j]_q!,
 \qquad D_j=(-q)^j,\qquad \gamma=Qq^{-d^2/2}.
\]
The coefficient of \(E^k\otimes F^k\) in the defining finite
\(R\) has the factorial ratio
\(([a]_q!/[a-k]_q!)^2/[k]_q!\).
Substitute \(p_j=(1-Q)^jq^{j(j-1)/2}[j]_q!\), together with its
Cartan factor. The resulting entry is exactly
\eqref{de:triangular-Fourier}, giving
\begin{equation}\label{de:factorial-similarity}
 A=\gamma(DG\otimes G)R
                 (G^{-1}\otimes G^{-1}D^{-1}).
\end{equation}
For example, at the entry in
\eqref{de:triangular-Fourier} the multiplier is
\(\gamma D_{a-k}G_{a-k}G_{b+k}/(G_aG_bD_b)\).
This fixes the positions of the factorial gauges; no estimate
on their condition numbers has been used.

The Fourier calculation involves a transposed crossing. To
compare it with the paper's \(B=\mathsf P R\), regard the
root self-duality as the antidiagonal matrix
\[
 Jv_j=c_jv_{d-j},\qquad
 c_j=(-1)^jq^{j(j-N)}.
\]
It obeys \(JKJ^{-1}=K^{-1}\) and
\begin{equation}\label{de:transpose-bridge}
 (J\otimes J)R(J^{-1}\otimes J^{-1})=R^{\mathsf T}.
\end{equation}
For an entry shifted from \((a,b)\) to \((a+k,b-k)\),
reflection leaves the factorial ratio unchanged because
\([N-j]_q=[j]_q\). Its \(J\)-ratio is
\[
 \frac{c_{d-a-k}c_{d-b+k}}{c_{d-a}c_{d-b}}
                 =q^{2k(a-b)+2k^2}.
\]
The reflected Cartan exponent is
\(H_aH_b/2+k(H_a-H_b)-2k^2\), with
\(H_a-H_b=2(b-a)\); the extra exponents cancel. This proves
\eqref{de:transpose-bridge} entry by entry, with the same
endpoint zeros.

For a braid on \(s\) strands put
\[
 U_s=(\mathsf F^{-1})^{\otimes s}
             \bigotimes_{j=1}^s(G^{-1}D^{j-1}),
 \qquad W_s=U_sJ^{\otimes s}.
\]
Transposing \eqref{de:factorial-similarity} gives the adjacent
crossing identity for \(\mathsf P R^{\mathsf T}\).
This crossing conserves the integer total, so it commutes with
\(D\otimes D\). The powers of \(D\) therefore combine into
the displayed single \(U_s\), valid at every adjacent position.
For a signed braid word \(\beta\) of writhe \(w\), one obtains
\begin{equation}\label{de:simultaneous-braid-gauge}
 \rho_K(\beta)=\gamma^w W_s\rho_B(\beta)W_s^{-1},
 \qquad
 W_sK^{\otimes s}W_s^{-1}=\mu_K^{\otimes s}.
\end{equation}
For the enhancement identity, \(G,D\) commute with \(K^{-1}\),
\(q^{-d}=-q\), and
\(\mathsf F^{-1}\operatorname{diag}(Q^j)\mathsf F=S\).
Negative generators are handled by inverting the same
adjacent-generator identity.

We give the cut-scalar argument explicitly. For an intertwiner
\(T\) of a tensor prefix with one last factor, set
\(C=\operatorname{Tr}_{\rm last}((I\otimes K)T)\).
The ordinary last trace of
\([T,E_{\rm prefix}\otimes K+I\otimes E]=0\)
gives \([C,E_{\rm prefix}]=0\). The \(K\)-weighted last trace
of the corresponding \(F\) equation gives
\([C,F_{\rm prefix}]=0\): its other two terms cancel by
cyclicity, using
\[
 T(K_{\rm prefix}^{-1}\otimes I)
 =(K_{\rm prefix}^{-1}\otimes K^{-1})T(I\otimes K).
\]
Commutation with the prefix \(K\) is immediate. The argument
can be iterated. Once only \(V_d\) remains, its distinct \(K\)
weights and nonzero adjacent \(F\) arrows force the cut
operator to be scalar.

In particular the enhanced partial traces of the crossing are
\begin{equation}\label{de:enhancement-traces}
 \operatorname{Tr}_2(B^{\pm1}(I\otimes K))
       =\alpha^{\pm1}I,\qquad
 \alpha=q^{(N^2-1)/2},\qquad \gamma\alpha=-q=\alpha_K.
\end{equation}
For the positive trace, evaluate on \(v_0\): only the input
pair \((0,0)\) contributes. For the inverse, evaluate on
\(v_d\): only \((d,d)\) contributes. Their values are
\(q^{d^2/2+d}\) and its inverse, which proves the formula.
These inverse calculations can also be read from
\[
 R^{-1}=
 \left(\sum_{k=0}^d
   \frac{(-1)^kq^{-k(k-1)/2}(q-q^{-1})^k}{[k]_q!}
                           E^k\otimes F^k\right)
                             q^{-H\otimes H/2}.
\]
The finite binomial identity makes its two polynomial factors
multiply to one below degree \(N\), and
\((E\otimes F)^N=0\).

Now take the \(K\)-weighted traces of every strand of a braid
closure except the chosen first one. Its cut operator \(C_B\)
is scalar by the preceding argument, irrespective of the number
of link components. In \eqref{de:simultaneous-braid-gauge},
all conjugations on traced factors cancel by cyclicity.
The remaining conjugation cancels because \(C_B\) is scalar.
Equation \eqref{de:enhancement-traces} cancels every writhe
factor, so
\begin{equation}\label{de:reduced-comparison}
                    \alpha_K^{-w}C_K=\alpha^{-w}C_B.
\end{equation}
Thus the comparison is exact even when a factorial gauge is
large, and it holds at even and composite levels.

Finally change from the \(K\)-pivot to the pivot of the marked
source. At completeness
\[
                   g=(-1)^dK,\qquad
 \theta=(-1)^d\alpha=q^{-d^2/2},\qquad |\theta|=1.
\]
For an \(s\)-strand closure the \(s-1\) traced factors contribute
\((-1)^{d(s-1)}\), and the curl normalization contributes
\((-1)^{-dw}\). If the closure has \(c\) components, its braid
permutation gives \(w\equiv s-c\pmod2\). Combining these signs
with \eqref{de:reduced-comparison} proves
\eqref{de:enhanced-comparison}. The turns and marked scalar
are exactly those already computed in
\eqref{de:root-turns}--\eqref{de:mark-on-support}.
For the same framed tangle, its planar ribbon reading and its
braid reading evaluate the same specified ribbon invariant.
The actual reduced cut is independent of its exposed edge by
the preceding source proof.

For a knot, correcting an integer framing multiplies the raw
scalar by the corresponding power of the unit \(\theta\);
the component sign is one. This gives
\eqref{de:raw-invariant}. For links with separately prescribed
component framings, retain those integer twist corrections
rather than replacing them by the total writhe, which also
counts crossings between different components.
\end{proof}

The full one-cut modified convention has a further, separate
normalization. If its modified dimension is defined by
\[
 d_N(\alpha)=-q^{N(1-N)/2}
 \frac{q^{\alpha+1}-q^{-\alpha-1}}
      {q^{N\alpha}-q^{-N\alpha}},
\]
then \(M=-q q^\alpha=e^u\), with the lift
\(\alpha=d+u/\log q\), gives the whole germ
\begin{equation}\label{de:modified-dimension}
 \begin{aligned}
 d_N\!\left(d+\frac{u}{\log q}\right)
   &=i^{N+3}\frac{\sinh u}{\sinh Nu},\\
 \lim_{u\to0}d_N\!\left(d+\frac{u}{\log q}\right)
   &=\frac{i^{N+3}}N.
 \end{aligned}
\end{equation}
The value at \(u=0\) is removable. It is attached once, to the
selected cut component, and not to every auxiliary circle.
This alternative convention is not used to change the
normalization in \eqref{de:raw-invariant}.
